% chapters/ch04_entropy.tex
\chapter{Entropy and Information Theory}
\label{ch:entropy}

\epigraph{%
  Information is not what you know. It is what you did not know before.%
}{\textit{---Flyxion}}

This chapter develops the information-theoretic language that frames the entire reconstruction enterprise. The central object is Shannon entropy, which measures the uncertainty of a probability distribution and therefore quantifies the difficulty of the inference problem at any stage of the algorithm~\cite{Cover2006}. We begin with the entropy of the uniform distribution over tree topologies---the initial uncertainty before any observations are made---and derive how each triplet observation reduces this entropy. The key monotonicity result, Proposition~\ref{prop:entropy_descent}, shows that entropy can only decrease as constraints accumulate: observations never increase uncertainty, though some may be redundant. We then derive a lower bound on the total number of queries needed for reconstruction, establishing that the algorithm's query complexity is fundamentally information-theoretically optimal up to logarithmic factors. The chapter closes with mutual information and KL divergence, tools that will reappear in the variational reformulation of tree inference in Appendix~\ref{app:variational}.

%% ── Figure: entropy descent ─────────────────────────────────
\begin{figure}[ht]
\centering
\begin{tikzpicture}
\begin{axis}[
  width=10cm, height=5.5cm,
  xlabel={Observation index $t$},
  ylabel={Residual entropy $H_t$ (bits)},
  xmin=0, xmax=50,
  ymin=0, ymax=42,
  xtick={0,10,20,30,40,50},
  grid=major, grid style={draw=nodegray!30},
  legend pos=north east,
  legend style={font=\small},
  thick
]
% Ideal descent: H_t = H_0 * (2/3)^t, H_0 = n log n ≈ 40 for n=10
\addplot[blue!70!black, thick, domain=0:50, samples=51]
  {40 * (0.585)^(x/3)};
\addlegendentry{ideal: $H_t = H_0 \cdot (2/3)^t$}

% More realistic: slower descent with plateaus
\addplot[red!70!black, thick, mark=*, mark size=1.2pt,
  mark options={fill=red!70!black}, only marks]
  coordinates {
    (0,40) (2,38.5) (4,36.8) (7,34.2) (10,31.5)
    (14,27.8) (18,23.4) (23,18.2) (29,12.5)
    (36,6.8) (42,2.9) (48,0.8) (50,0)
  };
\addlegendentry{typical oracle run}

\draw[dashed, gray!60] (axis cs:0,0) -- (axis cs:50,0)
  node[right, font=\footnotesize] {$H=0$};
\node[font=\footnotesize, nodeblue] at (axis cs:30,22)
  {reconstruction complete};
\end{axis}
\end{tikzpicture}
\caption{Entropy of the hypothesis space $H_t = \log|\mathcal{T}_n^{\mathcal{R}_t}|$ as a function of the number of triplet observations. Each observation reduces $H_t$ monotonically. The ideal curve assumes each triplet eliminates exactly $2/3$ of remaining trees; the dotted points show a typical oracle run with variable information gain per query.}
\label{fig:entropy_descent}
\end{figure}

\section{Shannon Entropy}

\begin{definition}[Shannon entropy]
Let $X$ be a discrete random variable on a finite set $\mathcal{X}$ with mass function $P(x)$. The \emph{Shannon entropy} is
\[
H(X) = -\sum_{x \in \mathcal{X}} P(x) \log P(x),
\]
with convention $0 \log 0 = 0$. Logarithms are in base 2 unless noted.
\end{definition}

\begin{proposition}[Bounds on entropy]
For $|\mathcal{X}| = k$: (1) $H(X) \geq 0$, with equality iff $X$ is deterministic; (2) $H(X) \leq \log k$, with equality iff $X$ is uniform.
\end{proposition}

\section{Joint Entropy, Conditional Entropy, and Mutual Information}

\begin{definition}[Conditional entropy and mutual information]
\[
H(X \mid Y) = H(X, Y) - H(Y), \qquad I(X; Y) = H(X) - H(X \mid Y).
\]
\end{definition}

\begin{proposition}
$I(X;Y) \geq 0$, with equality iff $X$ and $Y$ are independent.
\end{proposition}

\section{Entropy as Uncertainty over Tree Topologies}

Let $T$ be uniform over $\mathcal{T}_n$. Then
\[
H_0 := H(T) = \log |\mathcal{T}_n| \sim n \log n + O(n)
\]
by Corollary~\ref{cor:tree_entropy}. After observing $\mathcal{R}_t = \{r_1,\ldots,r_t\}$, the residual uncertainty is
\[
H_t = \log |\mathcal{T}_n^{\mathcal{R}_t}|.
\]

\begin{proposition}[Monotone entropy descent~{\cite{Cover2006}}]
\label{prop:entropy_descent}
$H_{t+1} \leq H_t$ for all $t \geq 0$.
\end{proposition}

\begin{proposition}[Information per triplet]
\label{prop:info_per_triplet}
Each consistent triplet observation eliminates approximately two-thirds of remaining trees, contributing approximately $\log(3/2) \approx 0.585$ bits of information under the uniform prior.
\end{proposition}

\section{Information-Theoretic Lower Bounds on Query Complexity}

\begin{theorem}[Query complexity lower bound]
\label{thm:query_lower_bound}
Any algorithm that reconstructs the tree topology with high probability must make at least $\Omega(n \log n)$ queries to a noiseless triplet oracle.
\end{theorem}

\begin{proof}
The total information to determine is $H_0 = \Theta(n \log n)$ bits. Each query reveals at most $\log 3 < 2$ bits. Therefore at least $H_0 / \log 3 = \Omega(n \log n)$ queries are required.
\end{proof}

\section{KL Divergence and Model Comparison}

\begin{definition}[KL divergence]
\[
\mathrm{KL}(P \| Q) = \sum_x P(x) \log \frac{P(x)}{Q(x)}.
\]
\end{definition}

$\mathrm{KL}(P\|Q) \geq 0$ with equality iff $P = Q$. KL divergence appears in the variational reformulation of tree inference (Appendix~\ref{app:variational}), where the inference objective decomposes into an expected log-likelihood and a KL regularization term.

\section{Exercises}

\begin{enumerate}
  \item Compute $H(T)$ for $n = 4, 5, 6$ under the uniform distribution over rooted binary tree topologies.
  \item Show that observing a consistent triplet reduces $H(T)$ by exactly $\log(3/2)$ bits when the prior is uniform and the three triplet relations are equally probable a priori.
  \item Show that the total number of triplet queries sufficient for noiseless reconstruction is $O(n \log n)$ by arguing each query can be chosen to eliminate at least a constant fraction of remaining consistent trees.
  \item Compute $\mathrm{KL}(P\|Q)$ when $P = \operatorname{Ber}(p)$ and $Q = \operatorname{Ber}(q)$, and show it is zero iff $p = q$.
\end{enumerate}
